\documentclass[11pt,a4paper]{article}
\usepackage[margin=1.8cm]{geometry}
\usepackage{graphicx}
\usepackage{booktabs}
\usepackage{amsmath,amssymb}
\usepackage{subcaption}
\usepackage{hyperref}
\usepackage{microtype}
\usepackage{siunitx}
\usepackage{xcolor}

\hypersetup{
    colorlinks=true,
    linkcolor=blue,
    citecolor=blue,
    urlcolor=blue
}

\title{\textbf{FLAG-Standard Systematic Error Budget for the QCD Chiral Transition Parameters along the LCP}}
\author{Junxiong Nie}
\date{\today}

\begin{document}
\maketitle

\begin{abstract}
Following the quality criteria and reporting conventions established by the Flavour Lattice Averaging Group (FLAG), this report presents the systematic error budget for the pseudo-critical chiral transition temperature $T_{pc}$ and the peak renormalized susceptibility $\chi_{\mathrm{peak}}/Z_m^2$ along the line of constant physics (LCP). All determinations are physically anchored at our most continuum-like, largest-volume lattice ensemble, $64^3 \times 16$ ($q_3$). The uncertainty decomposition incorporates statistical fluctuations, finite lattice discretization effects ($\Delta_{\mathrm{disc}}$), and finite volume effects ($\Delta_{\mathrm{FV}}$).
\end{abstract}

\section{Methodological Framework Following FLAG Conventions}

In high-precision lattice QCD determinations, the reporting of physical observables requires an unambiguous decomposition of statistical and systematic error sources.

\subsection{Physical Baseline Anchor ($64^3 \times 16$ Ensemble)}
In accordance with FLAG evaluation protocols, we adopt the state-of-the-art simulation on the $64^3 \times 16$ lattice ($q_3$) as the primary physical reference anchor:
\begin{equation}
    T_{pc}^{\mathrm{anchor}} = 155.12 \pm 0.18\text{ MeV}, \qquad
    \left(\chi_{\mathrm{peak}}/Z_m^2\right)^{\mathrm{anchor}} = 1.038 \pm 0.025\text{ GeV}^2.
\end{equation}
This ensemble possesses the smallest lattice spacing ($a \approx 0.08\text{ fm}$ at $N_t=16$) and the largest spatial box length ($L \approx 5.1\text{ fm}$ with aspect ratio $L/N_t = 4$), minimizing both ultraviolet discretization artifacts and infrared finite-volume distortions.

\subsection{Origin of Uncertainties and Independence of Resampling Samples}
\begin{enumerate}
    \item \textbf{Statistical Uncertainty ($\sigma_{\mathrm{stat}}$)}:
    The statistical error is purely derived from the Jackknife resampling covariance of the baseline $q_3$ ensemble across its $N = 1171$ independent Jackknife bins:
    \begin{equation}
        \sigma_{\mathrm{stat}} = \sqrt{\frac{N-1}{N} \sum_{k=1}^N \left( X_k(q_3) - \bar{X}(q_3) \right)^2}.
    \end{equation}
    This reflects the genuine statistical precision of the physical anchor.
    \item \textbf{Why Artificial Sample Alignment is Avoided}:
    Different lattice ensembles ($p_2, p_3, q_2, q_3$) originate from completely separate, uncorrelated Markov Chain Monte Carlo trajectories. The $k$-th Jackknife block in one ensemble shares no physical correlation with the $k$-th block of another ensemble. Truncating or aligning the sample count across different ensembles serves no statistical purpose and discards valid configurations. Each ensemble must retain its full, unadulterated sample set ($p_2: 2003, p_3: 1923, q_2: 2000, q_3: 1171$) to evaluate true ensemble distributions.
    \item \textbf{Finite Discretization Uncertainty ($\Delta_{\mathrm{disc}}$)}:
    Discretization effects arise from non-zero lattice spacing ($a \sim 1/N_t$). We isolate this effect by comparing the coarse lattice $p_3$ ($N_t=12$) with the fine baseline $q_3$ ($N_t=16$) at a matched large aspect ratio $L/N_t = 4$. To guarantee a conservative upper bound, the systematic uncertainty is evaluated as the maximum sample envelope deviation relative to the baseline anchor:
    \begin{equation}
        \Delta_{\mathrm{disc}} = \max_{k \in \mathrm{samples}(p_3)} |X_k(p_3) - \bar{X}(q_3)|.
    \end{equation}
    \item \textbf{Finite Volume Uncertainty ($\Delta_{\mathrm{FV}}$)}:
    Finite volume effects arise from infrared spatial box boundaries. We isolate this effect at fixed fine lattice spacing $N_t=16$ by comparing the smaller volume $q_2$ ($L=48$) against the large baseline volume $q_3$ ($L=64$). The systematic uncertainty is conservatively bounded by the maximum sample envelope deviation:
    \begin{equation}
        \Delta_{\mathrm{FV}} = \max_{k \in \mathrm{samples}(q_2)} |X_k(q_2) - \bar{X}(q_3)|.
    \end{equation}
\end{enumerate}

\section{Lattice Ensemble Inputs and Goodness-of-Fit}

Table~\ref{tab:flag_ensemble_inputs} compiles the geometric configurations, primary Padé fit orders, reduced chi-squared values $\chi^2/\mathrm{dof}$, and total valid Jackknife samples for the four core ensembles.

% Dynamically included LaTeX table generated by python
% ==============================================================================
% Auto-generated FLAG-Style Error Budget LaTeX Tables
% Generated by scripts/globalfit/error_budget_analysis.py
% ==============================================================================

\begin{table}[htbp]
    \centering
    \caption{Summary of lattice geometries, primary Padé fit orders, and full Monte Carlo Jackknife statistics across the four core LCP ensembles.}
    \label{tab:flag_ensemble_inputs}
    \small
    \begin{tabular}{lcccccc}
        \toprule
        \textbf{Ensemble} & \textbf{Lattice Dimensions} & \textbf{Padé Order} & $\boldsymbol{T_{pc}}$ [MeV] & $\boldsymbol{\chi_{\mathrm{peak}}/Z_m^2}$ [$\text{GeV}^2$] & $\boldsymbol{\chi^2/\mathrm{dof}}$ & \textbf{Jackknife Samples} \\
        \midrule
        $p_2$ & $36^3 \times 12$ & [2/2] & $152.74 \pm 0.89$ & $0.956 \pm 0.025$ & 0.775 & 2003 \\
        $p_3$ & $48^3 \times 12$ & [0/2] & $153.48 \pm 0.57$ & $0.935 \pm 0.018$ & 1.000 & 1923 \\
        $q_2$ & $48^3 \times 16$ & [2/2] & $154.06 \pm 0.75$ & $0.944 \pm 0.019$ & 0.840 & 2000 \\
        $q_3$ (Reference Anchor) & $64^3 \times 16$ & [0/2] & $155.12 \pm 0.18$ & $1.038 \pm 0.025$ & 1.952 & 1171 \\
        \bottomrule
    \end{tabular}
\end{table}

\begin{table}[htbp]
    \centering
    \caption{\textbf{FLAG-Style Error Budget Table}: Systematic uncertainty decomposition for the transition temperature $T_{pc}$ and peak susceptibility $\chi_{\mathrm{peak}}/Z_m^2$ anchored at the $64^3 \times 16$ ($q_3$) physical baseline.}
    \label{tab:flag_error_budget}
    \small
    \begin{tabular}{lcccc}
        \toprule
        \textbf{Uncertainty Source} & \multicolumn{2}{c}{$\boldsymbol{T_{pc}}$} & \multicolumn{2}{c}{$\boldsymbol{\chi_{\mathrm{peak}}/Z_m^2}$} \\
        \cmidrule(lr){2-3} \cmidrule(lr){4-5}
        & \textbf{Absolute [MeV]} & \textbf{Relative [\%]} & \textbf{Absolute [$\text{GeV}^2$]} & \textbf{Relative [\%]} \\
        \midrule
        \textbf{Central Reference Value ($64^3\times 16$)} & \textbf{155.12} & --- & \textbf{1.038} & --- \\
        \midrule
        Discretization / $a^2\to 0$ ($\Delta_{\mathrm{disc}}$) & $\pm 1.72$ & 1.11\% & $\pm 0.109$ & 10.47\% \\
        Finite Volume / $L\to\infty$ ($\Delta_{\mathrm{FV}}$) & $\pm 1.16$ & 0.75\% & $\pm 0.097$ & 9.35\% \\
        \midrule
        Combined Systematic ($\sigma_{\mathrm{sys}}$) & $\pm 2.07$ & 1.34\% & $\pm 0.146$ & 14.04\% \\
        \midrule
        \multicolumn{5}{l}{\textbf{Standard FLAG Result Presentation:}} \\
        \multicolumn{5}{c}{$T_{pc} = \mathbf{155.12 \pm 1.72_{\mathrm{disc}} \pm 1.16_{\mathrm{FV}}\text{ MeV}} = \mathbf{155.12 \pm 2.07_{\mathrm{sys}}} = \mathbf{155.12(207)\text{ MeV}}$} \\
        \multicolumn{5}{c}{$\chi_{\mathrm{peak}}/Z_m^2 = \mathbf{1.038 \pm 0.109_{\mathrm{disc}} \pm 0.097_{\mathrm{FV}}\text{ GeV}^2} = \mathbf{1.038 \pm 0.146_{\mathrm{sys}}} = \mathbf{1.038(146)\text{ GeV}^2}$} \\
        \bottomrule
    \end{tabular}
\end{table}

\section{FLAG Error Budget Table and Analysis}

Table~\ref{tab:flag_error_budget} presents the finalized FLAG-style systematic error budget, tabulating both absolute uncertainties and relative percentages for each physical observable.

\subsection{Physical Implications}
\begin{itemize}
    \item \textbf{Dominance of Discretization Artifacts}:
    For the pseudo-critical temperature $T_{pc}$, the finite discretization uncertainty $\Delta_{\mathrm{disc}} = 1.71\text{ MeV}$ ($1.10\%$) is the principal systematic component, while the finite volume uncertainty $\Delta_{\mathrm{FV}} = 1.16\text{ MeV}$ ($0.75\%$) is comparatively smaller due to the generous spatial extent ($L/N_t \ge 3$).
    \item \textbf{Total Systematic Precision}:
    Combining systematic components in quadrature yields $\sigma_{\mathrm{sys}} = 2.07\text{ MeV}$ ($1.33\%$) for $T_{pc}$, and $\sigma_{\mathrm{sys}} = 0.145\text{ GeV}^2$ ($13.98\%$) for the peak susceptibility $\chi_{\mathrm{peak}}/Z_m^2$.
    \item \textbf{FLAG Compact Presentation}:
    Expressing the final determinations in the standard FLAG concise bracket notation:
    \begin{align}
        T_{pc} &= \mathbf{155.12(18)_{\mathrm{stat}}(1.71)_{\mathrm{disc}}(1.16)_{\mathrm{FV}}\text{ MeV}} = \mathbf{155.12(18)(207)\text{ MeV}}, \\
        \chi_{\mathrm{peak}}/Z_m^2 &= \mathbf{1.038(25)_{\mathrm{stat}}(107)_{\mathrm{disc}}(97)_{\mathrm{FV}}\text{ GeV}^2} = \mathbf{1.038(25)(145)\text{ GeV}^2}.
    \end{align}
\end{itemize}

\end{document}
